\chapter{Probability, confidence, and calibration}
\label{ch:calibration}

This is the conceptual center of the project. A model can be accurate yet
\emph{miscalibrated}: its stated probabilities do not match how often it is right. Gating
actions on an untrustworthy probability is worse than not gating at all.

\section{Calibration, defined}

A model is \term{calibrated} if, among all the times it says ``0.8'', it is right about
80\% of the time. This is a property of \emph{groups} of predictions, not any single one.
TypeSafe states calibration exactly this way and warns it ``is not a guarantee about any
single answer.''\DOC{}

\begin{figure}[h]\centering
\begin{tikzpicture}
\begin{axis}[width=9cm,height=6cm,xlabel={stated probability},ylabel={observed accuracy},
  xmin=0,xmax=1,ymin=0,ymax=1,legend pos=north west,legend cell align=left,
  grid=both, grid style={clExternal!20}]
  \addplot[very thick, clExternal, dashed] coordinates {(0,0)(1,1)};
  \addlegendentry{perfectly calibrated}
  \addplot[very thick, clSecurity, mark=*, mark size=1.2] coordinates
    {(0.1,0.06)(0.3,0.16)(0.5,0.30)(0.7,0.45)(0.9,0.62)(1.0,0.72)};
  \addlegendentry{overconfident (typical)}
  \addplot[very thick, clAccent, mark=square*, mark size=1.2] coordinates
    {(0.1,0.20)(0.3,0.42)(0.5,0.58)(0.7,0.76)(0.9,0.93)(1.0,0.99)};
  \addlegendentry{underconfident}
\end{axis}
\end{tikzpicture}
\caption{A reliability diagram. Points on the dashed line are calibrated; below it is
overconfident, above it is underconfident.\datalabel{Illustrative (synthetic data)}}
\label{fig:reliability}
\end{figure}

\section{Measuring miscalibration: ECE and the noise floor}

The \term{Expected Calibration Error (ECE)} bins predictions by stated probability and
averages the gap between confidence and accuracy in each bin, weighted by bin size:
\[
\mathrm{ECE} = \sum_{b=1}^{B} \frac{n_b}{N}\,\bigl|\,\mathrm{acc}(b) - \mathrm{conf}(b)\,\bigr|.
\]
A subtlety that honest evaluations respect: with finite samples, even a perfectly calibrated
model shows a small nonzero ECE. That is the \term{noise floor}. Reporting ECE without its
floor overstates miscalibration.

\begin{figure}[h]\centering
\begin{tikzpicture}
\begin{axis}[width=11cm,height=5.2cm,ybar,bar width=10pt,
  symbolic x coords={OpenBookQA,CommonsenseQA,HellaSwag,OOD tickets},
  xtick=data, x tick label style={font=\scriptsize},
  ylabel={ECE}, ymin=0, ymax=0.13, legend pos=north west, legend cell align=left,
  enlarge x limits=0.15, grid=both, grid style={clExternal!20}]
  \addplot[fill=clModel!50, draw=clModel] coordinates
    {(OpenBookQA,0.024)(CommonsenseQA,0.032)(HellaSwag,0.029)(OOD tickets,0.107)};
  \addlegendentry{measured ECE}
  \addplot[fill=clExternal!30, draw=clExternal] coordinates
    {(OpenBookQA,0.024)(CommonsenseQA,0.019)(HellaSwag,0.018)(OOD tickets,0.024)};
  \addlegendentry{noise floor}
\end{axis}
\end{tikzpicture}
\caption{Jev's ECE is near the noise floor in distribution but about 4.4\(\times\) the floor
out of distribution.\datalabel{Reported by jev-ood-calibration}~\citep{jev_ood_calibration}}
\label{fig:ece}
\end{figure}

\section{Risk--coverage: the operating decision}

Gating is a trade: act only when confident (high \term{coverage} costs accuracy; low
coverage buys it). A \term{risk--coverage curve} shows accuracy on the automated fraction as
you sweep the threshold. The product's G0 gate asks for coverage $\ge 50\%$ at a chosen
precision target.

\begin{figure}[h]\centering
\begin{tikzpicture}
\begin{axis}[width=9cm,height=5.4cm,xlabel={coverage (fraction automated)},
  ylabel={accuracy on automated set}, xmin=0,xmax=1,ymin=0.6,ymax=1.02,
  grid=both, grid style={clExternal!20}]
  \addplot[very thick, clBackend, mark=*, mark size=1] coordinates
    {(0.1,0.99)(0.2,0.98)(0.3,0.97)(0.4,0.95)(0.5,0.93)(0.6,0.90)(0.8,0.84)(1.0,0.78)};
  \draw[clSecurity, dashed, thick] (axis cs:0.5,0.6) -- (axis cs:0.5,0.93);
  \node[clSecurity, font=\scriptsize, anchor=west] at (axis cs:0.52,0.66) {G0: coverage $\ge$ 0.5};
\end{axis}
\end{tikzpicture}
\caption{A risk--coverage curve: acting only on the most confident decisions raises
accuracy on the automated fraction.\datalabel{Illustrative (synthetic data)}}
\label{fig:riskcov}
\end{figure}

\section{Fixing calibration after the fact}

Two standard post-hoc recalibrators, fitted on a held-out labeled split:

\begin{description}[leftmargin=1.4cm,style=nextline]
  \item[Temperature scaling] divide the logits by a single learned scalar $T$ before the
    softmax~\citep{guo_calibration}. Simple, but it needs \emph{logits}. If the API returns
    only probabilities, or a probability is exactly 1.0, temperature scaling is undefined
    (\(\log 1 = 0\) after the softmax, \(\log\) of a saturated prob is problematic). This is
    a real constraint for Jev (see the pitfall).
  \item[Isotonic regression] fit a monotonic step function from stated probability to
    observed frequency. Works on probabilities directly and handled Jev well in an
    independent study, reaching ECE $\approx 0.01$ with a few hundred labels~\citep{jev_calibration_anthus}.
\end{description}

\begin{figure}[h]\centering
\begin{tikzpicture}
\begin{axis}[width=9cm,height=5.4cm,xlabel={stated probability},ylabel={observed accuracy},
  xmin=0,xmax=1,ymin=0,ymax=1,legend pos=north west,legend cell align=left,
  grid=both, grid style={clExternal!20}]
  \addplot[very thick, clExternal, dashed] coordinates {(0,0)(1,1)};
  \addlegendentry{ideal}
  \addplot[very thick, clSecurity, mark=*, mark size=1] coordinates
    {(0.1,0.06)(0.3,0.16)(0.5,0.30)(0.7,0.45)(0.9,0.62)(1.0,0.72)};
  \addlegendentry{before (raw)}
  \addplot[very thick, clBackend, mark=square*, mark size=1] coordinates
    {(0.1,0.11)(0.3,0.29)(0.5,0.49)(0.7,0.69)(0.9,0.88)(1.0,0.97)};
  \addlegendentry{after isotonic}
\end{axis}
\end{tikzpicture}
\caption{Post-hoc recalibration pulls an overconfident model toward the diagonal.\datalabel{Illustrative (synthetic data), effect direction reported by Jev-Calibration}~\citep{jev_calibration_anthus}}
\label{fig:recal}
\end{figure}

\begin{pitfall}
Independent testers found Jev returns probability \emph{exactly} 1.0 on many items, some of
them wrong~\citep{jev_baselines_eval}. A saturated 1.0 breaks temperature scaling and
collapses isotonic bins. The SDK's rule: verify whether logits are available; if not, use
isotonic or histogram binning with epsilon-clipping, and add a ``saturated'' band that
\emph{always escalates} for any authority above a plain filter. ``1.0 is not proof.''
\end{pitfall}

\begin{whyhere}
Independent studies also found the vendor \texttt{confidence} field was \emph{less}
calibrated than the maximum option probability~\citep{jev_ood_calibration}. So the SDK
defaults to gating on max probability, logs both, and lets the G0 study pick the winner.
Thresholds are fitted per decision point, per primitive, per model version, and refit when
wording, catalog, or version changes.
\end{whyhere}

\section{Conformal prediction (a brief note)}

\term{Conformal prediction} builds a set of options guaranteed (under exchangeability) to
contain the truth with a chosen probability, using a calibration set. Robotics work
(KnowNo) used it to decide when an agent should ask for help~\citep{knowno}. It is a natural
future tool for the escalation policy; the MVP starts with simpler threshold gating.

\begin{checkbox}
A model scores 92\% accuracy and ECE 0.15. Your policy acts whenever confidence
$\ge 0.9$. Why might this be dangerous, and which figure in this chapter would you inspect
first? (Appendix~\ref{app:answers}.)
\end{checkbox}
